\cotefile{157-5}
\scanpage{165}
\mspage{\unsure{81|87}}
% archive sheet number: 164
% margin note written slanted
\margin{\deleted{\unsure{Effacer|Effacé}}}

% lines 1-9
Une telle catégorie s'appellera \uline{grande devant $\pi$}.

(Il est immédiat que dans le cas où $J$ est elle-même ordonnée, cette terminologie est \unsure{cohérente} avec celle déjà adoptée pour les ens. ordonnés, i.e. que s'il existe une \unsure{application} croissante cofinale $\varphi:I\longrightarrow J$, avec $I$ grand devant $\pi$, alors $J$ est lui-$=$ grand devant $\pi$.) Une telle catégorie est automatiquement filtrante (Chap I \unsure{2.1 B b}).
% margin note written slanted
\margin{\unsure{Notons aussi} que \unsure{si on a} un foncteur cofinal $J'\longrightarrow J$, si $J'$ est grand devant $\pi$, $J$ l'est aussi (cp. \unsure{transitivité} de cofinalité).}

% lines 10-12
\uline{Proposition} 4.4.4. Soient $J$ une $\unsure{\mathcal{U}|\mathscr{U}}$-catégorie, \deleted{\unsure{filtrante}} $\pi$ un cardinal infini. Conditions équivalentes:

% lines 13-16
a) $J$ est grande devant $\pi$, i.e. il existe un ens. ordonné grand devant $\pi$ et un foncteur cofinal $\varphi:I\longrightarrow J$.

% lines 17-26
b) $J$ satisfait les deux conditions suivantes

$\mathrm{PF}_{\pi}1$: $\operatorname{Ob}J$ est grand devant $\pi$, i.e. pour toute famille $(y_{\alpha})_{\alpha\in A}$ d'objets de $J$, avec $\operatorname{card}A\leq\pi$, il existe un majorant $y$ dans $J$ (i.e. un $y$ et des flèches $y_{\alpha}\xrightarrow{v_{\alpha}}y$)

$\mathrm{PF}_{\pi}2$: Pour $y_0,y_1$ dans $J$ et \deleted{$\unsure{\forall} A$} \deleted{\unsure{pour tout} $\unsure{A}$} \deleted{deux} $(u_{\alpha})_{\alpha\in A}$ des flèches $y_0\xrightarrow{u_{\alpha}}y_1$, \inserted{\unsure{avec} $\operatorname{card}A\leq\pi$}, existe une flèche $y_1\xrightarrow{u'}y_2$ qui les égalise, i.e. telle que $u'u_{\alpha}=u'u_{\beta}$ $\forall\alpha,\beta\in A$.

(Il suffit de l'exiger pour $\operatorname{card}A\geq2$)

% margin note written slanted
\margin{Autres conditions. e) Pour $\forall$ ensemble $A$ avec $\operatorname{card}A\leq\pi$, $\delta_J:J\longrightarrow J^A$ est cofinal. e') Si $A_0$ est un ens. donné, de cardinal $=\pi$, $J\longrightarrow J^{A_0}$ est cofinal.}
% margin note written slanted
\margin{Voir aussi \unsure{4.9.6.9.1|4.4.6.9.1} et \unsure{4.9.6.9.2|4.4.6.9.2} pages \unsure{237,238}.}

% lines 27-32
c) Pour toute \deleted{\unsure{ensemble de flèches de $J$}} \inserted{partie $A$ de $\operatorname{Fl}J$}, avec $\operatorname{card}A\leq\pi$, existe une sous-catégorie $J'$ de $J$, avec $\operatorname{card}J'\leq\pi$, telle que $A\subset\operatorname{Fl}J'$ et que $J'$ admette un objet final.

% lines 33-34; sentence continues on the following scan
d) Soit $I$ l'ens des sous-catégories \inserted{$J'$} de $J$ telles que $\operatorname{card}J'\leq\pi$ et \unsure{qu'}$J'$ a un objet final,
